\subsection{R-22: arbitrary staircase size reaches the four-variable dimension ceiling}

\paragraph{Statement R-22.}
Let $b\geq3$ and $m\geq8b$ be integers, and define
\[
 T=m-4b+3,\qquad h=\left\lfloor\frac{m-6b+3}{b}\right\rfloor,
 \qquad J_{m,b}=(2b-2)hT.
\]
Then the quaternary permanent projection variety satisfies
\begin{equation}\label{R22:rank}
 c_m:=\dim\mathcal P_{m,4}\ \geq\ 2m^2+2+J_{m,b}.
\end{equation}
The proof uses one simultaneous deformation of the cyclic pencil.
Its untwisted motif parameters lie in the integer grid
$\{0,\ldots,4(b-1)T\}$; the translated copies have the prescribed
root-of-unity twists. A grid choice is allowed to depend on $m$ and on
the chosen primitive $m$th root of unity. No ranks from
different pencils are added.

In particular, for every $m\geq99$,
\begin{equation}\label{R22:asymptotic}
 4m^2-20m^{3/2}+2\ \leq\ c_m\ \leq\ 4m^2-2m+2.
\end{equation}
Thus $c_m=4m^2-O(m^{3/2})$. The upper bound is the row--column torus
bound proved in R-13, and an equality formula for $c_m$ is not asserted.

For any integers $n>m$ such that
\begin{equation}\label{R22:padding}
 2n^2\leq2m^2+J_{m,b}+2,
\end{equation}
the padded image has dimension greater than the determinant image.
Consequently R-14/G-10 give exactly-four-row multiplicity gaps in every
coordinate degree beyond the explicit G-10 threshold. The specified
stretched labels and effective bounds of G-17 also apply.
For every fixed real $1<\rho<\sqrt2$, this includes every integer
$m<n\leq\rho m$ as soon as
\begin{equation}\label{R22:ratio_threshold}
 m\geq\max\left\{99,
       \left\lceil\frac{100}{(2-\rho^2)^2}\right\rceil\right\}.
\end{equation}
In particular the eventual rectangular boundary of R-20 holds already
for every $m\geq2025$: rectangular gaps exist for every integer
$m<n<4m/3$, and G-12 excludes them for $n\geq4m/3$.
The additional range above $4m/3$ consists of nonrectangular gaps.

\begin{proof}
We give a uniform local calculation, prove its two Fourier rank
conditions, and then arrange that one integer choice works for all
interior bidegrees of the given matrix size.

\paragraph{1. The typed staircase and its simultaneous supports.}
Fix any primitive $m$th root $q$, and write
\[
 \ell_i=x-q^iy,\qquad v_i=z-q^iw.
\]
The cyclic pencil has entries $\ell_i$ at $(i,i)$ and $v_i$ at
$(i,i+1)$, with indices modulo $m$.
A block of $b$ rows, initially rows $0,\ldots,b-1$, has columns
$0,\ldots,b$. In addition to its diagonal and superdiagonal, allow
every entry $(i,j)$ except entries in column zero and the entry
$(0,b)$. Thus the permitted additional positions are
\[
 0\leq i<b,\quad1\leq j\leq b,\quad
 j\ne i,i+1,\quad(i,j)\ne(0,b).
\]
There are $b(b-2)$ such positions. Set their entries to
$\varepsilon C_{ij}$, where
\begin{equation}\label{R22:types}
 C_{ij}=\begin{cases}
 u_{ij}x+w_{ij}y,&j>i+1,\\
 u_{ij}z+w_{ij}w,&j<i.
 \end{cases}
\end{equation}
The symbol $w_{ij}$ denotes a scalar coefficient, distinct from the
coordinate $w$. All $2b(b-2)$ coefficients are initially independent
parameters, collectively denoted $\theta$.

Place copies of this same block at starts $a=bj$, $0\leq j<h$,
replacing $(y,w)$ in its additional forms by $(q^ay,q^aw)$.
All block rows are disjoint, since $bh\leq m-6b+3<m$.
For each block, its bottom $b-1$ rows have support only in its
$b$ columns numbered $1,\ldots,b$ relative to the start.
Its middle $b-1$ columns have support only in its $b$ rows numbered
$0,\ldots,b-1$. A neighboring block shares only a boundary column,
so both support assertions remain true with all blocks present.

These supports give exact cofactor factorizations. If an external row
is deleted, and one of the $b$ supported columns is deleted, all
bottom $b-1$ supported rows must occupy the remaining $b-1$ supported
columns in every matching. The complementary permanent is the product
of their local permanent minor and a common outside permanent.
The transposed assertion holds when an external column and one of
the $b$ supported rows are deleted. Other blocks can affect the outside
factor but cannot alter either factorization.

\paragraph{2. The second-order local defect for arbitrary size.}
Put $r=b-1$. Consider an $r\times(r+1)$ matrix $B(\varepsilon)$
with diagonal entries $a_i$, superdiagonal entries $b_i$, and
other entries $\varepsilon c_{ij}$, with $0\leq i<r$.
Form $\widetilde B$ by signing its additional entries with
\[
 \sigma_{ij}=\begin{cases}
 (-1)^{j-i-1},&j>i+1,\\
 (-1)^{i-j},&j<i.
 \end{cases}
\]
Let $P_j$ be the permanent minor deleting column $j$, and $D_j$
the corresponding determinant minor of $\widetilde B$.
There is one order-zero matching. A matching with one additional
edge has its other edges forced along the bidiagonal path; its
permutation sign is precisely the indicated $\sigma_{ij}$.
Hence
\begin{equation}\label{R22:local_defect}
 P_j-D_j=\varepsilon^2E_j+O(\varepsilon^3).
\end{equation}

We need the endpoint identities $E_0=E_r=0$ for all $r$.
For the minor deleting column $r$, a matching with no additional
forward edge is an upper-Hessenberg permutation. Each nontrivial
cycle is a consecutive interval, whose superdiagonal edges end in
one backward edge. Its permutation sign equals the sign assigned
to that backward edge. Two forward additional edges and no backward
edge cannot form a permutation. If there is exactly one forward
edge $i\mapsto j$, $j>i+1$, and one backward edge $h\mapsto l$,
there is a single nontrivial cycle, necessarily with
$l\leq i<j\leq h$. It follows superdiagonal edges from $l$ to $i$,
jumps to $j$, follows them to $h$, and returns to $l$.
Its length is $(i-l)+(h-j)+2$. The difference between its sign
exponent and the sum of the two assigned exponents is even, since
\[
 (j-i-1)+(h-l)-\big((i-l)+(h-j)+1\big)=2(j-i-1).
\]
Thus the endpoint minor agrees through order two.
Deleting column zero and transposing gives the same square
bidiagonal situation, with diagonal $b_i$ and superdiagonal
$a_{i+1}$. Under this operation a forward edge becomes a backward
edge with the same assigned sign, and conversely. This proves
\begin{equation}\label{R22:endpoints}
 E_0=E_r=0.
\end{equation}

Assume now that $a_i$ and the additional forward entries are
linear in $(x,y)$, while $b_i$ and the additional backward entries
are linear in $(z,w)$. The highest $(z,w)$ degree of $E_j$ is
$r-1$. Indeed, if there is no additional forward edge, the same
upper-Hessenberg argument gives equality of permanent and signed
determinant for every minor: the rows before the deleted column
form the upper-Hessenberg block, and all later superdiagonal
entries are forced. Therefore terms with no $(x,y)$ factor, and
terms with a diagonal $(x,y)$ factor and two backward edges, cancel.
At $(x,y)$ degree one the only remaining possibility uses one
forward edge, one backward edge, and otherwise superdiagonal edges.

For the minor deleting column $j$, compare its target columns with
the all-superdiagonal target set $\{1,\ldots,r\}$.
The target zero must come from the backward edge. Comparing the
two removed targets with the two added targets then forces the
additional edges to be
\[
 (j-1,h+1),\qquad(h,0),\qquad j\leq h\leq r-1.
\]
The minor permutation has the cycle
$(0,1,\ldots,j-1,h)$, of sign $(-1)^j$; the product of assigned
edge signs has opposite sign. Consequently
\begin{equation}\label{R22:Etop}
 [E_j]_{\deg_{z,w}=r-1}
 =2\sum_{h=j}^{r-1}c_{j-1,h+1}c_{h,0}
           \prod_{\substack{0\leq t<r\\t\ne j-1,h}}b_t
 \quad(1\leq j\leq r-1).
\end{equation}
This proves the formula for every block size, without extrapolating
a finite list of local minors.

The determinant row relation yields
\begin{equation}\label{R22:relation}
 \sum_{j=0}^{r}(-1)^{i+j}\widetilde B_{ij}P_j
 =\varepsilon^2(a_iE_i-b_iE_{i+1})+O(\varepsilon^3).
\end{equation}
For $i<r-1$, the highest component, divided by two, is
\[
 -\sum_{h=i+1}^{r-1}c_{i,h+1}c_{h,0}
                          \prod_{t\ne h}b_t,
\]
of bidegree $(1,r)$. For $i=r-1$ it is
\[
 a_{r-1}c_{r-2,r}c_{r-1,0}\prod_{t=0}^{r-3}b_t,
\]
of bidegree $(2,r-1)$.

In the original $b$-row block, denote these row forms by
$R_0,\ldots,R_{r-1}$. Explicitly,
\begin{align}
 R_i&=-\sum_{h=i+1}^{r-1}C_{i+1,h+2}C_{h+1,1}
                 \prod_{\substack{0\leq t<r\\t\ne h}}v_{t+1}
                     &&(0\leq i<r-1),\label{R22:Rforms}\\
 R_{r-1}&=\ell_r C_{r-1,r+1}C_{r,1}
                                      \prod_{t=1}^{r-2}v_t.\notag
\end{align}
The transposed supported-column calculation gives forms
$S_0,\ldots,S_{r-1}$:
\begin{align}
 S_0&=-\ell_1 C_{0,2}C_{r,1}\prod_{u=2}^{r-1}v_u,
                                                     \label{R22:Sforms}\\
 S_i&=\sum_{p=0}^{i-1}C_{p,i+1}C_{r,p+1}
                    \prod_{\substack{0\leq u<r\\u\ne p}}v_u
                    &&(1\leq i<r).\notag
\end{align}
Their $(x,y)$ degrees are respectively
\begin{equation}\label{R22:grades}
 (1,\ldots,1,2)\quad\hbox{and}\quad(2,1,\ldots,1).
\end{equation}
All their coefficients are homogeneous quadratic polynomials in
$\theta$. These displayed forms are highest components; the actual
limits from \eqref{R22:relation} can also have smaller $(z,w)$ degree.
We retain those lower components until the final independence argument.

\paragraph{3. The Fourier matrices and their exact meanings.}
Fix
\[
 2b-1\leq k,t\leq m-2b+1,\qquad K=q^k,\quad X=q^t.
\]
At $\varepsilon=0$, the common outside factors for the block at zero
are exactly
\begin{equation}\label{R22:outside}
 T_s^R=\ell_0\prod_{i=b}^{s-1}v_i\prod_{i=s+1}^{m-1}\ell_i,
 \qquad
 T_c^C=\prod_{i=b}^{c-1}\ell_i\prod_{i=c}^{m-1}v_i.
\end{equation}
The unique cyclic matching proves these formulas. For a local form
of bidegree $(a,g)$, $a+g=b$, take
\begin{equation}\label{R22:alignment}
 s=k+b-g=k+a,\qquad c=m+g-k.
\end{equation}
These indices are external to the corresponding supported rows or
columns throughout the stated range. The resulting leading products
have bidegree $(m-k,k)$.

Set $A(x,y)=\prod_{u=k}^{m-1}\ell_u$ and let $A_i$ be its coefficient
of $x^{m-k-i}y^i$. All $A_i$ are nonzero: the finite geometric-product
formula for these coefficients has numerator and denominator factors
$1-q^j$ with $0<j<m$ (the same formula used in R-16).
In cyclic character $t$, divide the coefficient with $y$-degree $i$
and $w$-degree $j$, $i+j=t$, by $A_i$.
Under this normalization the entire cyclic base image $W_0$ has the
two-dimensional binary coefficient space
\begin{equation}\label{R22:base}
 \operatorname{span}\left\{
 z\prod_{u=0}^{k-2}v_u,
 w\prod_{u=0}^{k-2}v_u\right\}.
\end{equation}
This is the base-character description proved in R-16; the full base
image has dimension $2m^2+2$ by R-15/R-17.

For a row form $F$ of $(x,y)$ degree $a\in\{1,2\}$, write
\begin{align}
 F&=\sum_{j=1}^{a+1}Q^R_{a,j}(x,y)C_j(z,w),\notag\\
 Q^R_{a,j}&=\prod_{u=j}^{a}(x-Kq^uy)
                       \prod_{u=1}^{j-1}(x-q^uy),\label{R22:rowbasis}\\
 \mathcal R_b(F)&=\sum_{j=1}^{a+1}X^j C_j(z,q^{-j}w)
       \prod_{u=b-j}^{b-2}v_u
       \prod_{u=-1}^{a-1-j}(z-Kq^uw).\label{R22:rowtransform}
\end{align}
For a column form use
\begin{align}
 F&=\sum_{j=0}^{a}Q^C_{a,j}(x,y)C_j(z,w),\notag\\
 Q^C_{a,j}&=\prod_{u=b-j}^{b-1}(x-q^uy)
               \prod_{u=0}^{a-1-j}(x-K^{-1}q^{b-a+u}y),
                                                     \label{R22:colbasis}\\
 \mathcal S_b(F)&=\sum_{j=0}^{a}X^{-j}C_j(z,Kq^{j-b}w)
       \prod_{u=-a+j}^{-1}v_u
       \prod_{u=0}^{j-1}(z-Kq^{-b+u}w).\label{R22:coltransform}
\end{align}
An empty product is one. All displayed bases are invertible for the
specified $q,k$. For the row bases, evaluation at $x=qy$, and then
at $x=q^2y$ when $a=2$, gives triangular independence; the possible
nonconstant vanishing factors are $q-1$, $K-1$, and $Kq-1$.
For the column bases evaluate at $x=q^{b-1}y$, then
$x=q^{b-2}y$. The same conditions, up to nonzero powers, suffice.
None can vanish in the interior range. Thus these transformations are
defined over $\mathbb Q(q)$ after $K=q^k$ and introduce no denominators
depending on $\theta$ or $X$.

Here are direct verifications of the Fourier completions.
Multiplying the row outside $(x,y)$ factor by $Q^R_{a,j}$ completes
$A(x,q^jy)$. Coefficient division gives $q^{ji}=X^jq^{-j\,\deg_w}$.
The remaining $(z,w)$ roots have indices $b-j,\ldots,k+a-1-j$.
Their common factor is
\[
 H_R=\prod_{u=b-1}^{k-2}v_u,
\]
and the residual degree-$b$ form is \eqref{R22:rowtransform}.
For the column completion, multiplication by $Q^C_{a,j}$ completes
$A(x,q^{b-k-j}y)$. After division and the nonzero character scalar
$X^{k-b}$, it gives $X^{-j}q^{(k-b+j)\deg_w}$.
The outside $(z,w)$ roots become $-a+j,\ldots,k-b-1+j$.
Their common factor and residual form are
\[
 H_C=\prod_{u=0}^{k-b-1}v_u,
 \qquad \mathcal S_b(F).
\]
Thus both transformations represent the actual character quotients
of the cofactor limits, including all shifts in $k$ and $t$.

Let $M_R$ be the $(b+1)\times(b+1)$ coefficient matrix of
\begin{equation}\label{R22:MR}
 z\prod_{u=0}^{b-2}v_u,\quad
 w\prod_{u=0}^{b-2}v_u,\quad
 \mathcal R_b(R_0),\ldots,\mathcal R_b(R_{r-1}).
\end{equation}
Let $M_C$ be the corresponding coefficient matrix of
\begin{equation}\label{R22:MC}
 z\prod_{u=-b}^{-2}(z-Kq^uw),\quad
 w\prod_{u=-b}^{-2}(z-Kq^uw),\quad
 \mathcal S_b(S_0),\ldots,\mathcal S_b(S_{r-1}).
\end{equation}
These are precisely base plus row, and base plus column, after
removing their respective common factors. For the column matrix the
additional nonzero character scalar just used also rescales the base;
its root indices are then $k-b,\ldots,k-2$, as in
\eqref{R22:MC}.

\paragraph{4. Uniform row nonvanishing at every interior bidegree.}
Fix the primitive root $q$ and any interior $k$. We prove
$\det M_R\not\equiv0$ as a polynomial in $\theta,X$ by a specialization
that may use elements of $\mathbb Q(q)$.
In the local bottom-row matrix, retain only nearest forward chords
$c_{i,i+2}=p_i$ and first-column backward chords $c_{h,0}=z$.
Set
\[
 p_i=\ell_1\quad(0\leq i\leq r-3),\qquad
 p_{r-2}=\ell_2,
\]
and set its other chords to zero. Parameters outside this row matrix
can be arbitrary. Its row forms are
\[
 -p_i z\prod_{\substack{1\leq u\leq r\\u\ne i+2}}v_u
 \quad(0\leq i\leq r-2),
 \qquad \ell_r\ell_2z\prod_{u=1}^{r-2}v_u.
\]
For $i\leq r-3$ their normalized forms are
\begin{equation}\label{R22:diagonal_forms}
 -X^2z\prod_{\substack{-1\leq u\leq r-1\\u\ne i}}v_u.
\end{equation}
Evaluate modulo the base space at the $r$ distinct roots
$z/w=q^h$, $0\leq h\leq r-1$. The first $r-2$ columns in
\eqref{R22:diagonal_forms} are diagonal on the first $r-2$ roots and
zero on the last two; their diagonal entries are nonzero for $X\ne0$.

The remaining two local forms are
\[
 U=-\ell_2z\prod_{u=1}^{r-1}v_u,\qquad
 V=\ell_r\ell_2z\prod_{u=1}^{r-2}v_u.
\]
Write
\[
 \alpha=\frac{q-1}{K-1},\qquad
 \beta=\frac{K-q}{K-1},\qquad
 \ell_2=\alpha Q^R_{1,1}+\beta Q^R_{1,2}.
\]
Their normalized forms $P_U=\mathcal R_b(U)$,
$P_V=\mathcal R_b(V)$ are
\begin{align*}
 P_U={}&-z\left[
  \alpha X\prod_{u=0}^{r-2}v_u\,(z-Kq^{-1}w)
 +\beta X^2\prod_{u=-1}^{r-3}v_u\,v_{r-1}\right],\\
 P_V={}&c_1Xz\prod_{u=0}^{r-3}v_u\,(z-Kq^{-1}w)(z-Kw)\\
 &{}+c_2X^2z\prod_{u=-1}^{r-4}v_u\,v_{r-1}(z-Kq^{-1}w)\\
 &{}+c_3X^3z\prod_{u=-2}^{r-5}v_u\,v_{r-2}v_{r-1}.
\end{align*}
where $c_j$ are the coefficients of $\ell_r\ell_2$ in the basis
$Q^R_{2,j}$. Only the first is needed:
\begin{equation}\label{R22:c1}
 c_1=\frac{(1-q^{r-1})(1-q)}{(1-K)(1-Kq)}.
\end{equation}
It follows by evaluating the basis expansion at $x=qy$.

Put $A=q^{r-2}$, $B=q^{r-1}$ and evaluate at $w=1$.
$P_U(A)$ has no $X$ term; $P_U(B)$ and $P_V(B)$ have only their
$X$ terms. The coefficient of $X^2$ in the last two-by-two
evaluation determinant is therefore
\begin{equation}\label{R22:last_minor}
 \begin{split}
 [X^2]\big(P_U(A)P_V(B)-P_V(A)P_U(B)\big)
 ={}&\alpha c_1AB
 \prod_{u=0}^{r-2}(B-q^u)
 \prod_{u=0}^{r-3}(A-q^u)\\
 &\cdot(B-K/q)(A-K/q)(A-K).
 \end{split}
\end{equation}
Every factor is nonzero for the fixed primitive $q$ and every
$2b-1\leq k\leq m-2b+1$. In particular the potential equalities
$K=q^{b-1},q^{b-2},q^{b-3}$ lie outside the range, as do $K=1,q^{-1}$.
This also covers $b=3$, when the first diagonal block is empty.
The kernel of evaluation at these $r$ roots is exactly the
two-dimensional base space in \eqref{R22:MR}. Thus the row matrix
has nonzero determinant as a polynomial in $X$ at this specialization,
for every interior $k$ and every primitive $q$ under consideration.
Its specialized degree in $X$ is at most $2b-3$: the first $b-3$
evaluation columns contribute $X^{2(b-3)}$, and the last determinant
has degree at most three.

\paragraph{5. Column nonvanishing at the same primitive root and bidegree.}
The complete parameter domain in \eqref{R22:types} is invariant under
the reflection of additional positions
\begin{equation}\label{R22:reflection}
 (i,j)\longmapsto(b-j,b-i).
\end{equation}
It preserves $j-i$, hence preserves the forward $(x,y)$ type and
backward $(z,w)$ type. Transpose the global cyclic pencil and reverse
its indices by setting $M'_{uv}=M_{b-v,b-u}$ modulo $m$.
Its diagonal and cycle entries are
\[
 \ell_{b-u},\qquad v_{b-u-1}.
\]
With
\[
 q'=q^{-1},\qquad y'=q^by,\qquad w'=q^{b-1}w,
\]
they become the standard cyclic entries $x-(q')^uy'$ and
$z-(q')^uw'$. This is an invertible change of parameters and variables
over $\mathbb Q(q)$.

More explicitly, if $C$ is the transposed supported-column matrix of
the original block and $B'$ is the supported-row matrix of the new
block, then $B'_{ij}=C_{r-1-i,r-j}$. Reversing both rows and columns
preserves the determinant signs of each square minor. It changes the
row relation $a_iE_i-b_iE_{i+1}$ to the negative of the corresponding
reversed relation. This proves \eqref{R22:Sforms} from
\eqref{R22:Rforms} and identifies the two collections of actual
cofactor limits. An external original column $c=m+g-k$ becomes the
external new row
\[
 b-c\equiv k+b-g\pmod m,
\]
exactly the row alignment in \eqref{R22:alignment}. Bidegree $k$ is
unchanged. The change of $y,w$ only rescales character coordinates by
nonzero numbers; it does not mix them or change their available
coefficient interval.

For precision, the following argument does not assume that separate
generic row and column ranks are already simultaneous. Apply the
row specialization just proved at the primitive root $q^{-1}$ and
at this same $k$. Its determinant is a nonzero polynomial in $X$
of degree at most $2b-3$. There are $T=m-4b+3>2b-3$ interior
characters, so at least one of their distinct root values gives
nonzero determinant. Restricting the resulting binary forms to the
allowed coefficient interval preserves their independence: the
Toeplitz/Schur argument in the next paragraph applies to
$H_R\operatorname{Sym}^b$, using any $b+1$ consecutive allowed
indices. Reflection therefore gives a column specialization at $q,k$
with independent actual column classes. If $M_C$ were identically
singular as a function of its parameters and $X$, such independence
after restriction would be impossible. Hence
\begin{equation}\label{R22:two_nonzero}
 \det M_R(q,q^k,\theta,X)\not\equiv0,
 \qquad
 \det M_C(q,q^k,\theta,X)\not\equiv0
\end{equation}
for every interior $k$, over the same parameter ring and the same
fixed field $\mathbb Q(q)$.

\paragraph{6. Joint row--column rank and coefficient restriction.}
Suppose both matrices are nonsingular for some interior $k,t$.
Before coefficient restriction, base plus row spans
$H_R\operatorname{Sym}^b\mathbb C^2$, and base plus column spans
$H_C\operatorname{Sym}^b\mathbb C^2$. Their common factor is
\[
 H=\prod_{u=b-1}^{k-b-1}v_u,
 \qquad N:=\deg H=k-2b+1\geq0.
\]
After dividing by $H$, the two spaces are
\[
 \left(\prod_{u=k-b}^{k-2}v_u\right)\operatorname{Sym}^b\mathbb C^2,
 \qquad
 \left(\prod_{u=0}^{b-2}v_u\right)\operatorname{Sym}^b\mathbb C^2.
\]
The two degree-$b-1$ factors are coprime, since their indices are
distinct and lie below $m$. Their intersection is exactly their
product times $\operatorname{Sym}^1\mathbb C^2$, namely the base.
Their joint span is therefore all
\begin{equation}\label{R22:joint_space}
 H\operatorname{Sym}^{2b-1}\mathbb C^2,
\end{equation}
of dimension $2b$. This is the required joint independence, not an
addition of two ranks without checking their intersection.

The allowed $w$-degree interval in character $t$ is
\[
 \max\{0,t-(m-k)\}\ \leq j\leq\ \min\{k,t\}.
\]
Its length is at least $2b$, since $k,m-k,t,m-t\geq2b-1$.
Choose $2b$ consecutive indices $e,\ldots,e+2b-1$ in it;
then $0\leq e\leq N$. Restriction of multiplication by $H$ to
these indices has determinant
\begin{equation}\label{R22:schur}
 \det(H_{e+i-j})_{0\leq i,j<2b}
 =\pm s_{((2b)^e)}(q^{b-1},q^b,\ldots,q^{k-b-1})\ne0.
\end{equation}
The identity is dual Jacobi--Trudi, as in R-16/R-19 and R-B15.
For nonvanishing, there are $N$ geometric variables, and the Weyl
numerator exponents for this rectangle are distinct in
$[0,N+2b-1]=[0,k]$. Since $k<m$, their powers of $q$ are distinct,
and the numerator is a nonzero Vandermonde determinant; the
denominator is nonzero for the same reason. For $N=0$ the coefficient
matrix is the identity. This proves that all $2b-2$ normal classes
remain independent modulo the entire base character space.
The same proof with $b+1$ in place of $2b$ and the corresponding
geometric factor proves the row-only restriction fact used in the
preceding paragraph.

\paragraph{7. One integer grid choice for all bidegrees.}
We now fix $m,b,q$ and choose the motif parameters once.
By \eqref{R22:grades} and \eqref{R22:rowtransform}, the $X$ exponents
of the row determinant lie in
\[
 [b-1,2b-1]:
\]
every one of its $b-1$ normal columns has minimum exponent one,
the $b-2$ columns with $a=1$ have maximum two, and the column with
$a=2$ has maximum three. The column determinant has exponents in
$[-b,0]$, by \eqref{R22:coltransform} and the sum
$2+(b-2)=b$ of its $a$ values. Consequently, for each interior $k$,
\begin{equation}\label{R22:Fk}
 F_k(\theta,X)=X\det M_R(q,q^k,\theta,X)
                    \det M_C(q,q^k,\theta,X)
\end{equation}
is an ordinary polynomial in $X$ of degree at most $2b$.
Its coefficients lie in $\mathbb Q(q)[\theta]$, and its total
parameter degree is at most $4(b-1)$. By
\eqref{R22:two_nonzero} it is nonzero. The product over all $T$
interior values of $k$ is therefore a nonzero polynomial with
parameter degree at most $D=4(b-1)T$.

A polynomial of degree at most $D$ in each parameter over a field
of characteristic zero
cannot vanish at all points of $\{0,\ldots,D\}^{2b(b-2)}$ unless
it is zero: prove this successively in each variable by univariate
interpolation. Apply this over the field $\mathbb Q(q)(X)$.
There is an integer grid choice $\theta$ for which the product is
nonzero as a polynomial in $X$. Fix that one choice. Every
$F_k(\theta,X)$ is then nonzero, so at most $2b$ interior characters
are bad for each $k$. At least
\begin{equation}\label{R22:good_modes}
 T-2b=m-6b+3
\end{equation}
characters have all $2b-2$ independent normal classes.

This argument deliberately permits the grid choice to depend on
$m,q$; it makes no claim that one fixed integer motif works for
every root order. It needs neither a bound for $\varphi(m)$ nor a
generic root-of-unity specialization inferred from numerical tests.
If a concrete choice is wanted, enumerate the finite grid in
lexicographic order and test \eqref{R22:Fk} in exact cyclotomic
arithmetic, choosing the first with nonzero product. This is a finite
algorithm, with no polynomial-time assertion.

\paragraph{8. Translations, all bidegrees, and a single differential.}
For fixed $k$, the leading vectors from the block at start $a=bj$
are the transforms by $\tau^{bj}$ of those at zero, where
$\tau:(y,w)\mapsto(qy,qw)$. Suppose a linear combination of these
vectors lies in $W_0$. In a good character $t$, their independent
normal classes give $2b-2$ separate equations of the form
\[
 \sum_{j=0}^{h-1}c_j(q^{bt})^j=0.
\]
Every value of $q^{bt}$ has at most $\gcd(b,m)\leq b$ preimages.
The good characters in \eqref{R22:good_modes} therefore contain
at least $h$ distinct such values. Ordinary Vandermonde independence
forces every $c_j$ to vanish. Hence each of the $T$ bidegrees supplies
$(2b-2)h$ independent leading classes modulo $W_0$.

The actual limits of \eqref{R22:relation} need not be homogeneous
for the $(x,y)/(z,w)$ bigrading. In any relation among them modulo
$W_0$, take the largest selected $(z,w)$ degree. The corresponding
highest components are the independent classes just proved; all
remaining components have smaller degree. Since $W_0$ is bigraded,
these coefficients vanish. Descending in degree proves independence
of all $(2b-2)hT$ full limits modulo $W_0$.

Finally each cofactor combination in \eqref{R22:relation}, multiplied
by its common outside permanent, belongs to the differential image
of the same simultaneous pencil whenever $\varepsilon\ne0$.
Dividing by $2\varepsilon^2$ gives a polynomial family, by
\eqref{R22:local_defect}. Its limit is the full vector just used.
Choose $2m^2+2$ ordinary differential columns independent at
$\varepsilon=0$, and append these polynomial families. The resulting
matrix has a nonzero minor at zero, by the quotient independence.
That minor remains nonzero for a common generic nonzero
$\varepsilon$. The differential rank is therefore at least
$2m^2+2+(2b-2)hT$, proving \eqref{R22:rank}.

\paragraph{9. Quantitative consequences and the precise ceiling.}
Since $h\geq m/b-7$ and $T\geq m-4b$, with both lower bounds
nonnegative under $m\geq8b$,
\[
 J_{m,b}\geq(2b-2)(m/b-7)(m-4b)
 \geq(2-2/b)m^2-22bm.
\]
For $m\geq99$, take $b=\lfloor\sqrt{m/11}\rfloor$.
Then $b\geq3$, $m\geq8b$, and
$\frac12\sqrt{m/11}\leq b\leq\sqrt{m/11}$. Thus
\[
 \frac{2m^2}{b}+22bm\leq6\sqrt{11}\,m^{3/2}
 <20m^{3/2},
\]
which proves \eqref{R22:asymptotic}. The padded image has dimension
$c_m+3$ by R-14, while the determinant image has dimension
$2n^2+2$ by R-17. This gives \eqref{R22:padding} and all the stated
multiplicity consequences. Under \eqref{R22:ratio_threshold},
\[
 4m^2-20m^{3/2}+2\geq2\rho^2m^2+2\geq2n^2+2,
\]
so the dimension comparison is strict after padding. Substitution
$\rho=4/3$ gives the threshold $2025$ and, with G-11/G-12 or the
rectangular specialization of G-17, the claimed exact eventual
rectangular range.

For comparison, the dimension method itself cannot pass the
asymptotic ratio $\sqrt2$. Its necessary strict inequality, using
the upper bound in \eqref{R22:asymptotic}, is
\[
 2n^2<4m^2-2m+3.
\]
In particular it fails for $n\geq\sqrt2\,m$ and $m\geq2$.
The proved lower bound approaches this ceiling from below; for
example $m\geq169$ and
$m<n\leq\sqrt2\,m-4\sqrt m$ suffice by direct squaring.
This does not exclude other methods of finding multiplicity gaps
beyond the dimension comparison, and it does not address
superlinear padding or polynomial-time multiplicity computation.
\end{proof}

